
\subsection{The \texorpdfstring{$(\infty,2)$}{(infty,2)}-functoriality of traces}

Throughout this subsection, fix a symmetric monoidal $(\infty,2)$-category $(\cC,\otimes)$ with
monoidal unit $\1_{\cC} \in \cC$ and write $\Omega\cC$ for the $(\infty,1)$-category of
endomorphisms $\1_{\cC} \to \1_{\cC}$. The symmetric monoidal structure on $\cC$
induces a natural symmetric monoidal structure on $\Omega \cC$, such that
the tensor product $f\otimes_{\Omega\cC} g$ of two endomorphisms
$f,g\colon \1_{\cC} \to \1_{\cC}$ is simply given by the map $f\otimes g\colon \1_{\cC} \to \1_{\cC}$. 

\begin{defn}
  An object $X\in \cC$ is said to be \tdef{dualizable}
  if the functor $X\otimes - \colon \cC \to \cC$ admits an adjoint of the form
  $Y\otimes - \colon \cC \to \cC$. In this case, we write $Y= X^{\vee}$ and refer to $X^{\vee}$ as
  the \tdef{dual} of $X$. We denote by $\mdef{\cC^{\dual}}\subseteq \cC$ the full subcategory
  spanned by the dualizable objects.
\end{defn}

\begin{rmk}
  For any $X\in \cC^{\dual}$ we have a canonical equivalence $(X^{\vee})^{\vee}\simeq X$ and
  the assignment $X \mapsto X^{\vee}$ refines to an equivalence of categories
  $\cC^{\dual, \op}\rar{\sim} \cC^{\dual}$.
\end{rmk}

By definition, any $X\in \cC^{\dual}$ comes equipped with evaluation and
co-evaluation maps
\[ \coev_{X}\colon \1_{\cC} \too X\otimes X^{\vee}\]
\[ \ev_{X}\colon X^{\vee} \otimes X \too \1_{\cC}\]
Moreover, since $\cC$ is symmetric monoidal, there is a canonical \textit{braiding},
giving a natural equivalence
\[ \tau \colon X\otimes X^{\vee} \simeq X^{\vee} \otimes X \,.\]

\begin{defn}
  For any map $f\colon X\to X \in \cC^{\dual}$, the \tdef{trace} of $f$ is defined 
  to be the composite
  \[\mdef{\tr{f}{\cC}} \colon  \1_{\cC}  \xrightarrow{\coev_X}  X \otimes X^{\vee}
    \rar{f \otimes \id} X \otimes X^{\vee} \rar{\tau} X^{\vee} \otimes X 
    \rar{\ev_X} \1_{\cC}\,.\]
  We refer to $\tr{\id_X}{\cC}\colon \1_{\cC} \to \1_{\cC}$ as the \tdef{dimension} of
  $X$ and denote it by $\mathrm{dim}_{\cC}(X)$.
\end{defn}

% \begin{lem}\label{lem: sym mon preserve trace}
%     Let $F\colon \cC\to \cD$ be a symmetric monoidal $(\infty,2)$-functor
%     and let $f\colon Y\otimes X \to X\otimes Z$ be a generalized
%     endomorphism of some dualizable object $X\in \cC^{\dual}$.  
%     We have a natural equivalence 
%     \[ F(\tr{f}{\cC})\simeq \tr{F(f)}{\cD}\,, \]
%     in the $(\infty,1)$ category of maps $F(Y)\to F(Z)$. 
% \end{lem}
We will make use of the $(\infty,2)$-functoriality of the trace.
For this purpose, recall the following construction of 
Hoyois--Scherotzke--Sibilla~\cite{Hoyois_Scherotzke_Sibilla2017}.

\begin{cons}
The assignment $\cC\mapsto \cC^{\dual}$ refines to a functor
\[ (-)^{\dual}\colon \Cat^{\otimes}_{(\infty,2)} \too \Cat_{(\infty,2)}\]
which preserves filtered colimits and all limits by~\cite[Proposition 4.6.1.11]{Lurie2017}
and thus admits a left adjoint denoted by
\[ \rm{Fr}^{\dual}\colon \Cat_{(\infty,2)} \too \Cat^{\otimes}_{(\infty,2)}.\]
The $(\infty,1)$-category of traceable endomorphisms $\Endtrl(\cC)$ is defined as the
maximal sub-$(\infty,1)$-category of the $(\infty,2)$-category 
$\Fun^{\rm{oplax}}_{\otimes}(\rm{Fr}^{\dual}(B\NN), \cC)$
of symmetric monoidal oplax functors from the ``dualizably free'' $(\infty,2)$-category
on the walking category with one endomorphism.
\end{cons}

Explicitly, objects of $\Endtrl(\cC)$ are given by self
maps of dualizable objects $f\colon X \to X$ in $\cC$, while a morphism $(X,f) \to (Y,g)$
consists of an internal left adjoint $\varphi\colon X \to Y$ together with a 2-morphism
$\alpha \colon \varphi \circ f \to g \circ \varphi $ i.e.~a lax commutative diagram
\[\begin{tikzcd}
    X & Y \\
    X & Y \arrow["\varphi", from=1-1, to=1-2] \arrow["f"', from=1-1, to=2-1] \arrow["g", from=1-2, to=2-2] \arrow["\alpha", shorten <=4pt, shorten >=4pt, Rightarrow, from=2-1, to=1-2] \arrow["\varphi"', from=2-1, to=2-2]
  \end{tikzcd}\]
  Note that the assignment $\cC \mapsto \Endtrl(\cC)$ is manifestly
  functorial, meaning any symmetric monoidal $(\infty,2)$-functor
  $F\colon \cC \to \cD$ induces a functor $F\colon \Endtrl(\cC)\to \Endtrl(\cD)$. The following
  theorem says that that the trace is compatible with everything in sight.

\begin{thm}\cite[Theorem 1.7]{Hoyois_Scherotzke_Sibilla2017}\label{thm: trace functor}
  The $(\infty,1)$-category $\Endtrl(\cC)$ inherits a natural
  symmetric monoidal structure, such that the assignment taking $(X,f) \in \Endtrl(\cC)$ to the
  trace $\tr{f}{\cC} \in \Omega_{\cC}$ refines to a symmetric monoidal functor
  \[ \Trc_{\cC}\colon \Endtrl(\cC) \too \Omega \cC.\]
  Moreover, for any symmetric monoidal $(\infty,2)$-functor 
  $F\colon \cC \to \cD$, we have a commutative diagram:
  \[ \begin{tikzcd}
	{\Endtrl(\cC)} & {\Endtrl(\cD)} \\
	{\Omega \cC} & {\Omega \cD}
	\arrow["F", from=1-1, to=1-2]
	\arrow["{\Trc_{\cC}}"', from=1-1, to=2-1]
	\arrow["{\Trc_{\cD}}", from=1-2, to=2-2]
	\arrow["{\Omega F}"', from=2-1, to=2-2]
\end{tikzcd}\]
\end{thm}

\begin{rmk}
More explicitly, the functoriality of $\Trc_{\cC}$ works as follows:

Given a map $(\varphi, \alpha)\colon (X,f) \to (Y,g)$, denote by $\psi \colon Y\to X$ the internal right adjoint of $\varphi$.
Then $\Trc_{\cC}$ takes the map $(\varphi,\alpha)$ to the natural transformation given by the composite
\[ \tr{f}{\cC} \to \tr{\psi \varphi f}{\cC} \rar{\alpha} \tr{\psi g \varphi}{\cC}
\simeq \tr{\varphi \psi g}{\cC} \to \tr{g}{\cC},\]
where the first and last map are given by the unit and counit of the adjunction
$\varphi \dashv \psi$, respectively. 
\end{rmk}


\subsection{Span categories}
% Diagrammatically, these are also induced by the 2-cells
% in the following diagram
% \begin{equation}\label{eq: trace functoriality}
%   \begin{tikzcd}
% 	{\1_\cC} & {X \otimes X^\vee} && {X\otimes X^\vee} & {\1_\cC} \\
% 	{\1_\cC} & {Y\otimes Y^\vee} && {Y \otimes Y^\vee} & {\1_\cC}
% 	\arrow["{\rm{coev}_X}", from=1-1, to=1-2]
% 	\arrow[equals, from=1-1, to=2-1]
% 	\arrow["{f\otimes \id}", from=1-2, to=1-4]
% 	\arrow[Rightarrow, from=1-2, to=2-1]
% 	\arrow["{\varphi \otimes \psi^\vee}", from=1-2, to=2-2]
% 	\arrow["{\ev_X}", from=1-4, to=1-5]
% 	\arrow["\alpha"', Rightarrow, from=1-4, to=2-2]
% 	\arrow["{\varphi \otimes \psi^\vee}"', from=1-4, to=2-4]
% 	\arrow[Rightarrow, from=1-5, to=2-4]
% 	\arrow[equals, from=1-5, to=2-5]
% 	\arrow["{\rm{coev}_Y}"', from=2-1, to=2-2]
% 	\arrow["{g\otimes \id}"', from=2-2, to=2-4]
% 	\arrow["{\ev_Y}"', from=2-4, to=2-5]
% \end{tikzcd}\end{equation}

% % \begin{rmk}\label{rmk:compatibility-of-fancy-trace-with-normal-trace}
% %   The functor $\Trc_{\cC}$ enjoys the following compatibility with the natural symmetric
% %   monoidal structure on $\Omega\cC$.
% %   Given an endomorphism $\gamma \colon a \to a$ of dualizable object in $\Omega \cC$, we
% %   obtain an endomorphism in $\Endtrl(\cC)$ of the form
% %   \[ (a,\gamma) \colon (\1_{\cC},\id) \too (\1_{\cC},\id). \]
% %   Then we have a natural equivalence
% %   \[ \Trc(a,\gamma) \simeq \tr{\gamma}{\Omega \cC}\]
% %   of endomorphisms of
% %   $\Trc(\1_{\cC},\id)= \dim{\1_{\cC}}{\cC} = \1_{\Omega \cC} \in \Omega\cC$.\todo{Insert reference to Hoyois et al or short argument.}
% % \end{rmk}

% The $(-)^{\vee}\colon (\cC^{\dual})^{\op}\to \cC^{\dual}$ induces a functor on
% $\Endtrl(\cC^{\dual})$ in the obvious way. Note that in general $\Endtrl(\cC)^{\dual}\neq \Endtrl(\cC^{\dual})$.

% \begin{lem}
%   Let $F\colon \cC\to \cD$ be a unital lax symmetric monoidal functor of symmetric monoidal $\infty$-categories. Let $X\in \cC$ be an object such that $F(X)\in \cD^{\dual}$.
%   Given any $\phi \colon X^{\vee}\to Y \in \cC$ the induced map $F(\phi)\colon F(X^{\vee}) \to F(Y)$
%   factors naturally through the comparison map
%   \[\delta\colon F(X^{\vee})\to F(X)^{\vee}\]
%   given by the adjoint of the composite 
%   \[ F(X^{\vee})\otimes F(X) \to F(X^{\vee}\otimes X) \xrightarrow{F(\ev)} F(\1_{\cC})\simeq \SS.\]
% \end{lem}
%
\begin{defn}\label{defn: geometric setup}
A \tdef{geometric setup} is a pair $(\mathcal C, \mathcal E)$ where $\mathcal C$ is an $\infty$-category with admits all finite limits
and $\mathcal E \subseteq \Map (\mathcal C)$ is a class of morphisms
in $\cC$, satisfying the following conditions:
\begin{enumerate}[label=(\roman*)]
    \item $\mathcal E$ contains all equivalences and is closed under composition.
    \item $\mathcal E$ is stable under base change in $\mathcal C$.
\end{enumerate}
 Morphisms in $\mathcal E$ are called \tdef{admissible} and we also call
 an object $X\in \cC$ admissible whenever the terminal map $X\to \pt$ is.
Given a geometric setup $(\cC,\cE)$, we write $\mdef{\Span{\cC, \all, \cE}}$ for
the associated symmetric monoidal $\infty$-category of \tdef{spans}, see \todo{Add Ref}. 
If $\cE=\all$ is the class of all maps, we suppress it from the notation and
simply write $\Span{\cC}= \Span{\cC,\all,\all}$. 
\end{defn}

\begin{lem}\label{lem: duality datum}
    Let $(\cC,\cE)$ be a geometric context and suppose we have $X\in \Span{\cC, \all, \cE}$ 
    such that the terminal map $X\to \pt$ and the diagonal 
    $\Delta \colon X \to X\times X$ are admissible. 
    Then $X$ is dualizable and canonically self dual, with duality datum given by the 
    spans:
    \[  \coev_X\colon \left( \pt \xleftarrow{} X \xrightarrow{\Delta} X\times X\right)
      \quad \text{and}\quad  \ev_{X}\colon \left( X\times X \xleftarrow{\Delta} X \xrightarrow{} \pt
    \right)\,.\]
\end{lem}

\begin{lem}\label{lem: trace in Span}
   Let $(\cC,\cE)$ be a geometric context and let $X\in \cC$ such that the terminal
   map $X\to \pt$ and the diagonal $\Delta \colon X \to X\times X$ are admissible.
   For any morphism
   \[\Lambda= (X\xleftarrow{f}X\xrightarrow{g}X)\] 
   in $\Span{\cC, \all, \cE}$ we have that 
   \[\tr{\Lambda}{\Span{\cC, \all, \cE}} = (\pt \leftarrow X^{\Lambda} \to \pt ) \in \Span{\cC, \all, \cE}\,,\]
   where $X^{\Lambda}$ is defined to be the pullback: 
   \[\begin{tikzcd}
	{X^{\Lambda}} & X \\
	X & {X\times X\,.}
	\arrow[from=1-1, to=1-2]
	\arrow[from=1-1, to=2-1]
	\arrow["\lrcorner"{anchor=center, pos=0.125}, draw=none, from=1-1, to=2-2]
	\arrow["{(f,g)}", from=1-2, to=2-2]
	\arrow["\Delta", from=2-1, to=2-2]
\end{tikzcd}\]
\end{lem}
\begin{proof}
 This is straightforward from the definition of the trace and the duality datum described  
 in \cref{lem: duality datum}. 
\end{proof}
To fully leverage the functoriality of trace, we will also
make use of an $(\infty,2)$-categorical enhancement of the span
category.

\begin{defn}\label{defn: Span 2-cat}
Let $(\cC,\cE)$ be a geometric context. Given a pair of subclasses
$\cI, \cP \subseteq \cE$ 
such that $\cI$ and $\cP$ contain all equivalences, are closed
under composition, base change and diagonals we write
$\mdef{\sSpan{\cC, \all, \cE}_{\cP,\cI}}$ for the symmetric monoidal 
$(\infty,2)$-category of \cite[Definition 4.10]{Cnossen_Lenz_Linskens2026}. 
\end{defn}

\begin{rmk}
By construction, the  underlying $(\infty,1)$-category of $\sSpan{\cC, \all, \cE}_{\cP,\cI}$
agrees with $\Span{\cC, \all, \cE}$. Moreover, a 2-morphism from 
$\left(X\leftarrow Z\to Y\right)$ to $\left(X\leftarrow Z'\to Y\right)$ is given 
by a commutative diagram:
\[\begin{tikzcd}
	& Z & \\
	X & W & Y \\
	& Z^{\prime}
	\arrow[from=1-2, to=2-1]
	\arrow[from=1-2, to=2-3]
	\arrow["i"', from=2-2, to=1-2]
	\arrow["f", from=2-2, to=3-2]
	\arrow[from=3-2, to=2-1]
	\arrow[from=3-2, to=2-3]
\end{tikzcd}\]
where $i\in\cI$ and $f\in\cP$.
\end{rmk}

\begin{rmk}\label{rmk: 2span and cohomology}
 The $(\infty,2)$-category of spans is useful for the following: Suppose we are
 given a six-functor formalism, i.e.~a lax symmetric monoidal $(\infty,2)$ functor
 \[ \cD\colon \sSpan{\cC, \all, \cE}_{\cP,\cI}\too \cV\]
 where $\cV$ is some presentably monoidal $(\infty,2)$ category. The $\infty$-category
 of endomorphisms of the terminal object $\pt \in \cC$ is given by 
 \[ \Omega_{\pt} \sSpan{\cC, \all, \cE}_{\cP,\cI} \simeq \Span{\cC, \cP, \cI}\]
 \todo{cite?}
Thus by passing to endomorphism $(\infty,1)$-categories, we obtain an associated cohomology
 theory 
 \[ \cH_{\cD}(-)\colon \Span{\cC,\cP,\cI} \too \Omega_{\cD(\pt)}(\cV)\,.\]
\end{rmk}

We record the following technical lemma about internal adjoints in the
$(\infty,2)$-category of spans:

\begin{lem}\label{lem: internal adjoints}
  Let $(\cC,\cE)$ be a geometric context and let $\cI,\cP\subseteq \cE$ be as in 
  \cref{defn: Span 2-cat} . For any admissible $X\in \cC$, the 
  following hold:
  \begin{enumerate}[label=(\alph*)]
    \item If the maps $X\to \pt$ and $\Delta\colon X\to X\times X$ lie in $\cP$ , 
      the 1-morphism \[(\pt\leftarrow X\xrightarrow\id X)\in \Span{\cC, \all ,\cE}_{\cP,\cI}\]
      has an internal right adjoint given by the flipped span.
  \item If the maps $X\to \pt$ and $\Delta\colon X\to X\times X$ lie in $\cI$, the 
    1-morphism \[(X\xleftarrow{\id} X\to \pt)\in \Span{\cC, \all, \cE}_{\cP,\cI}\] 
    has an internal right adjoint given by flipping the span.
  \end{enumerate}
\end{lem}
\begin{proof}
In the case of (a), the adjunction between $\pt\leftarrow X\xrightarrow\id X$ and $X\xleftarrow{\id} X\to \pt$ is witnessed by the 2-morphisms:
\[\begin{tikzcd}
	& X\times X & \\
	\pt & X\times X & \pt \\
	& \pt
	\arrow[from=1-2, to=2-1]
	\arrow[from=1-2, to=2-3]
	\arrow["\id"', from=2-2, to=1-2]
	\arrow[ from=2-2, to=3-2]
	\arrow[from=3-2, to=2-1]
	\arrow[from=3-2, to=2-3]
\end{tikzcd}
\quad \text{and} \quad 
\begin{tikzcd}
	& X & \\
	 X & X  & X \\
	& X\times X
	\arrow[from=1-2, to=2-1]
	\arrow[from=1-2, to=2-3]
	\arrow["\id"', from=2-2, to=1-2]
	\arrow["\Delta", from=2-2, to=3-2]
	\arrow[ from=3-2, to=2-1]
	\arrow[ from=3-2, to=2-3]
\end{tikzcd}\]
  As for (b), simply mirror the diagrams above along the horizontal axis.
\end{proof}

%%% Local Variables:
%%% TeX-master: "main"
%%% End:
